% SPDX-License-Identifier: Apache-2.0
%
% Abstract. The rule here is that every number is one we can point at a file
% for. No number in this abstract is aspirational.

Post-quantum cryptography is being deployed under a deadline. NIST finalised
ML-KEM and ML-DSA in 2024, the German BSI published its national migration
guidance in October 2026 requiring cryptographic inventories by the end of that
year, and the migration horizon for non-high-risk systems is 2035. What is
scarce is not algorithms. It is \emph{evidence}: a machine-checkable account of
what was proved, what was assumed, and what remains open.

We report an engineering case study of one PQC implementation, FIBEMATE, along
two axes that are usually reported separately. The first is formalization. We
mechanize the VWZ construction in Isabelle/HOL, tiering 56 obligations by what
each one depends on. Two tiers are complete with no remaining \texttt{sorry}
(T0, purely formal rewriting, 17 lemmas; T1-A, algebraic axioms, 16 lemmas), one
is nearly complete (T1-D, 5 lemmas with 2 open), and two are substantially open
(T2, linear algebra, 10 lemmas with 7 open; T3, security and probabilistic
obligations, 8 lemmas all open because their subjects are still undefined). The
repository holds 185 lemmas and 1 theorem in total, of which 17 carry
\texttt{sorry}.

The second axis is supply chain. We generate a CycloneDX 1.6 conformant
Cryptographic Bill of Materials covering 9 cryptographic components across 304
distinct source files (647 component--file occurrences), validated against the
official schema in CI.

The methodological contribution is what happened when we pointed verification
tools at our own artefacts. Three findings are reported in detail. First, the
Isabelle session runs with \texttt{quick\_and\_dirty}, so a successful build
proves only that the non-\texttt{sorry} parts type-check; the exit code is
therefore not a progress metric, and we argue that reporting one is a form of
misrepresentation. Second, our first CBOM failed official schema validation with
51 errors across 7 categories, none of which our own 18 internal tests detected,
because those tests validated the CBOM against our parser rather than against
the standard. Third, our public CBOM viewer had never rendered its table at all:
it threw on first call, and the defect was invisible to code review and to the
tests, and would have been reported as working had we not executed it.

We argue that these are not incidental bugs but the characteristic failure mode
of self-validating engineering, and that the discipline required to catch them
--- external references, negative controls, and execution rather than inspection
--- is a prerequisite for PQC evidence to be worth anything to an auditor.

\medskip
\noindent\textbf{Keywords:} post-quantum cryptography, formal verification,
Isabelle/HOL, CBOM, CycloneDX, software supply chain, BSI TR, reproducibility.
